% -------------------------------------------------------------
% METADATOS
% -------------------------------------------------------------

% --- DATOS DE LA PORTADA ---
\newcommand{\unefaNucleo}{NÚCLEO ARAGUA - SEDE MARACAY}
% UNIDAD (DESCOMENTAR SI EXISTE)
%\newcommand{\unefaUnidad}{E.A.D. DE INVESTIGACIÓN}
\newcommand{\carrera}{Escribir Carrera Aquí}
\newcommand{\tipoTrabajo}{Trabajo Especial de Grado}
\newcommand{\datosTitulo}{ANÁLISIS DE FALLAS EN SISTEMAS EMBEBIDOS UTILIZANDO REDES NEURONALES}
\newcommand{\tituloCorto}{COLOCAR ENCABEZADO CORTO AQUÍ}

% --- DATOS DEL AUTOR(ES) Y TUTOR ---
% TUTOR
\newcommand{\datosTutorNombre}{Federico Pinto}
\newcommand{\datosTutorCI}{V-1.234.567}
\newcommand{\datosTutorTitulo}{Ing.}
% PRIMER AUTOR
\newcommand{\datosAutorNombre}{Pepito Pérez}
\newcommand{\datosAutorCI}{V-9.876.543}
% SEGUNDO AUTOR (DESCOMENTAR SI EXISTE)
%\newcommand{\datosSegAutorNombre}{Petra Gómez}
\newcommand{\datosSegAutorCI}{V-10.987.654}
% TÍTULO DE AMBOS AUTORES
\newcommand{\datosAutoresTitulo}{Br.}

% --- DATOS DE LUGAR Y FECHA ---
% CIUDAD PRESENTACIÓN
\newcommand{\cuidadPresentacion}{Maracay}
% ESTADO PRESENTACIÓN
\newcommand{\estadoPresentacion}{Aragua}

% 1. El día de presentación (ej. 28)
\newcommand{\diaPresentacion}{28}

% 2. El MES de presentación se define como un NÚMERO (1=Enero, 9=Septiembre, 11=Noviembre)
\newcommand{\mesNumPresentacion}{11} % <--- ¡Define el mes como un número!

% 3. El año (ej. 2025)
\newcommand{\anoPresentacion}{2025}
